\documentclass[12pt]{article}
\usepackage[utf8]{inputenc}
\usepackage[T1]{fontenc}
\usepackage{hyperref}
\title{East Damar Language Resources for K-12 Curriculum Development}
\author{Compiled by Claude Sonnet 5.5 (claude-sonnet-5-5) for Andr\'as Kornai}
\date{}
\begin{document}
\maketitle

\section{Introduction}
East Damar (ISO 639-3 \texttt{dmr}, also called South Damar) is an Austronesian language of Damar Island, Maluku, Indonesia, with roughly 2,800 speakers cited in older surveys; other sources give higher figures. It is not closely related to West Damar, the other language of the island. The web search found very little material suitable for schooling: no primers, textbooks or school programmes were found. This report lists only what appeared in search results.

\section{Primary Educational Resources}
No primers, textbooks or curricula in East Damar were found in the search results.

\section{Linguistic and Reference Materials}
\begin{itemize}
\item East Damar language, Wikipedia overview: \url{https://en.wikipedia.org/wiki/East_Damar_language}
\item Academic paper on Damar Island (van Engelenhoven), Leiden repository: \url{https://scholarlypublications.universiteitleiden.nl/access/item%3A2883235/download}
\item Numerals entry, Max Planck Institute: \url{https://lingweb.eva.mpg.de/channumerals/Damar-East.htm}
\item ASJP word-list entry: \url{https://asjp.clld.org/languages/EAST_DAMAR}
\item Austronesian Comparative Dictionary language index: \url{https://ling.lll.hawaii.edu/dicts/ACD/acd-l_e.htm}
\item Verbix page: \url{https://docs.verbix.com/Languages/DamarEast}
\end{itemize}

\section{Religious and Community Resources}
\begin{itemize}
\item Joshua Project profile of the East Damar people: \url{https://joshuaproject.net/people_groups/11481}
\end{itemize}

\section{Digital and Audio Resources}
\begin{itemize}
\item Global Recordings Network audio and slideshow materials: \url{https://globalrecordings.net/en/language/dmr}
\end{itemize}

\section{Curriculum Development Context}
The language is reported as stable and not known to be taught in schools. Damar is a district of the regency of Maluku Barat Daya. Any curriculum would have to begin with new documentation and community-led materials, since no existing teaching resources surfaced. Speaker figures differ between sources and should be verified locally.
\end{document}